%! The Rational Parent Function \& Transformations — Unit 5, Lesson 5.4 — Slides (Beamer)
\documentclass[aspectratio=169, 11pt]{beamer}
\usepackage{algebra2-beamer}   % provides \forestheader{} and \sectionlabel[color]{}
% Standard coordinate grid + axes: {xmin}{xmax}{ymin}{ymax}
\newcommand{\lgrid}[4]{%
  \draw[step=1, linegray!45, thin] (#1,#3) grid (#2,#4);
  \draw[->, thick, charcoal] (#1,0)--(#2,0) node[below right, font=\tiny]{$x$};
  \draw[->, thick, charcoal] (0,#3)--(0,#4) node[left, font=\tiny]{$y$};}

\begin{document}

% TITLE SLIDE (CourseName is not defined in beamer, so the name is literal)
{
\setbeamercolor{background canvas}{bg=forest}
\begin{frame}[plain]
  \vspace{2.8cm}\hspace{0.55cm}%
  \begin{minipage}{0.9\textwidth}
    {\color{goldacc}\bfseries\small LESSON 5.4}\\[0.5cm]
    {\color{white}\bfseries\LARGE The Rational Parent Function}\\[0.35cm]
    {\color{white}\bfseries\large and Transformations}\\[1.0cm]
    {\color{forestmid}\small Algebra 2: Shepherd~$\cdot$~Unit 5: Rational Functions}
  \end{minipage}
\end{frame}
}

% HOOK
\begin{frame}[t, plain]
  \forestheader{Same outputs. Three inputs later.}
  \sectionlabel{Hook}
  \begin{columns}[T]
    \begin{column}{0.52\textwidth}
      \centering
      \renewcommand{\arraystretch}{1.25}
      \begin{tabular}{|c|c||c|c|}
      \hline
      $x$ & $y=\frac1x$ & $x$ & $y=\frac{1}{x-3}$ \\ \hline
      $-2$ & $-\frac12$ & $1$ & $-\frac12$ \\ \hline
      $-1$ & $-1$       & $2$ & $-1$ \\ \hline
      $0$  & undef.     & $3$ & undef. \\ \hline
      $1$  & $1$        & $4$ & $1$ \\ \hline
      $2$  & $\frac12$  & $5$ & $\frac12$ \\ \hline
      \end{tabular}
    \end{column}
    \begin{column}{0.44\textwidth}
      \begin{block}{The sentence for the whole period}
        \textbf{Shifting a rational function carries its asymptotes with it.} \\[3pt]
        The blow-up moved from $x=0$ to $x=3$. The settling-down value stayed at $y=0$. \\[3pt]
        \textbf{The wall moved. The floor did not.}
      \end{block}
    \end{column}
  \end{columns}
\end{frame}

% THE PARENT
\begin{frame}[t, plain]
  \forestheader{The parent: $f(x)=\dfrac{1}{x}$}
  \sectionlabel{Definition}
  \begin{columns}[T]
    \begin{column}{0.42\textwidth}
      \centering
      \begin{tikzpicture}[scale=0.36]
        \lgrid{-5}{5}{-5}{5}
        \begin{scope} \clip (-5,-5) rectangle (5,5);
          \draw[very thick, forest, domain=-5:-0.2, samples=100, smooth] plot (\x,{1/(\x)});
          \draw[very thick, forest, domain=0.2:5,  samples=100, smooth] plot (\x,{1/(\x)});
        \end{scope}
      \end{tikzpicture}
    \end{column}
    \begin{column}{0.55\textwidth}
      \begin{itemize}\setlength{\itemsep}{3pt}
        \item A \textbf{hyperbola} --- two separate \textbf{branches}.
        \item VA $x=0$; \; HA $y=0$.
        \item Domain: all reals, $x\neq0$. \; Range: all reals, $y\neq0$.
        \item Decreasing on $(-\infty,0)$ \emph{and} on $(0,\infty)$ --- two intervals, never one.
        \item Origin symmetry: an \textbf{odd} function.
      \end{itemize}
      \vspace{3pt}
      {\footnotesize End behavior: $y\to0$ in both directions.}
    \end{column}
  \end{columns}
\end{frame}

% NO INTERCEPTS
\begin{frame}[t, plain]
  \forestheader{The fact nobody believes: this parent has \emph{no} intercepts}
  \sectionlabel[goldacc]{Hold the line}
  \begin{columns}[T]
    \begin{column}{0.48\textwidth}
      \begin{block}{No $x$-intercept}
        A fraction is zero only when its \textbf{numerator} is zero. \\[3pt]
        The numerator is $1$. It is never zero.
      \end{block}
    \end{column}
    \begin{column}{0.48\textwidth}
      \begin{block}{No $y$-intercept}
        A $y$-intercept is $f(0)$. \\[3pt]
        But $0$ is not in the domain --- there is nothing to plot.
      \end{block}
    \end{column}
  \end{columns}
  \vspace{8pt}
  {\footnotesize Every parent before this one passed through the origin: $y=x$, $y=x^2$, $y=|x|$,
  $y=x^3$. This one misses both axes entirely, and that is the first thing that makes it different.}
\end{frame}

% THE TWO SHIFTS
\begin{frame}[t, plain]
  \forestheader{Each shift moves one asymptote}
  \sectionlabel[goldacc]{The heart of the lesson}
  \begin{columns}[T]
    \begin{column}{0.48\textwidth}
      \centering
      \begin{tikzpicture}[scale=0.30]
        \lgrid{-4}{7}{-4}{4}
        \draw[navy, dashed, thick] (3,-4)--(3,4);
        \node[navy, font=\tiny, above] at (3,4) {$x=3$};
        \begin{scope} \clip (-4,-4) rectangle (7,4);
          \draw[very thick, forest!35, dashed, domain=-4:-0.25, samples=90, smooth] plot (\x,{1/(\x)});
          \draw[very thick, forest!35, dashed, domain=0.25:7, samples=90, smooth] plot (\x,{1/(\x)});
          \draw[very thick, forest, domain=-4:2.75, samples=90, smooth] plot (\x,{1/((\x)-3)});
          \draw[very thick, forest, domain=3.25:7, samples=90, smooth] plot (\x,{1/((\x)-3)});
        \end{scope}
      \end{tikzpicture} \\[2pt]
      {\footnotesize $y=\dfrac{1}{x-3}$ --- right $3$, so the \textbf{VA} goes to $x=3$.}
    \end{column}
    \begin{column}{0.48\textwidth}
      \centering
      \begin{tikzpicture}[scale=0.30]
        \lgrid{-5}{5}{-3}{6}
        \draw[navy, dashed, thick] (-5,2)--(5,2);
        \node[navy, font=\tiny, below right] at (3.4,2) {$y=2$};
        \begin{scope} \clip (-5,-3) rectangle (5,6);
          \draw[very thick, forest!35, dashed, domain=-5:-0.33, samples=90, smooth] plot (\x,{1/(\x)});
          \draw[very thick, forest!35, dashed, domain=0.33:5, samples=90, smooth] plot (\x,{1/(\x)});
          \draw[very thick, forest, domain=-5:-0.2, samples=90, smooth] plot (\x,{1/(\x)+2});
          \draw[very thick, forest, domain=0.25:5, samples=90, smooth] plot (\x,{1/(\x)+2});
        \end{scope}
      \end{tikzpicture} \\[2pt]
      {\footnotesize $y=\dfrac{1}{x}+2$ --- up $2$, so the \textbf{HA} goes to $y=2$.}
    \end{column}
  \end{columns}
  \vspace{4pt}
  \begin{block}{}
    \centering\textbf{The asymptotes are the parent's axes, and they travel with it.} Draw the two
    dashed lines first and the branches have nowhere else to go.
  \end{block}
\end{frame}

% STRETCH AND FLIP
\begin{frame}[t, plain]
  \forestheader{$a$ stretches and flips --- and moves no asymptote}
  \sectionlabel{Scaling}
  \begin{columns}[T]
    \begin{column}{0.56\textwidth}
      \centering
      \begin{tikzpicture}[scale=0.26]
        \begin{scope}
          \lgrid{-5}{5}{-5}{5}
          \begin{scope} \clip (-5,-5) rectangle (5,5);
            \draw[very thick, forest!35, dashed, domain=-5:-0.2, samples=90, smooth] plot (\x,{1/(\x)});
            \draw[very thick, forest!35, dashed, domain=0.2:5, samples=90, smooth] plot (\x,{1/(\x)});
            \draw[very thick, forest, domain=-5:-0.6, samples=90, smooth] plot (\x,{3/(\x)});
            \draw[very thick, forest, domain=0.6:5, samples=90, smooth] plot (\x,{3/(\x)});
          \end{scope}
          \node[charcoal, font=\tiny] at (0,-6.6){$y=\frac3x$};
        \end{scope}
        \begin{scope}[xshift=13cm]
          \lgrid{-5}{5}{-5}{5}
          \begin{scope} \clip (-5,-5) rectangle (5,5);
            \draw[very thick, forest!35, dashed, domain=-5:-0.2, samples=90, smooth] plot (\x,{1/(\x)});
            \draw[very thick, forest!35, dashed, domain=0.2:5, samples=90, smooth] plot (\x,{1/(\x)});
            \draw[very thick, forest, domain=-5:-0.2, samples=90, smooth] plot (\x,{-1/(\x)});
            \draw[very thick, forest, domain=0.2:5, samples=90, smooth] plot (\x,{-1/(\x)});
          \end{scope}
          \node[charcoal, font=\tiny] at (0,-6.6){$y=-\frac1x$};
        \end{scope}
      \end{tikzpicture}
    \end{column}
    \begin{column}{0.42\textwidth}
      \begin{itemize}\setlength{\itemsep}{3pt}
        \item $|a|>1$: branches pushed \textbf{away} from the asymptotes.
        \item $0<|a|<1$: branches pulled \textbf{in}.
        \item $a<0$: branches jump to the \textbf{other pair} of regions.
      \end{itemize}
      \vspace{4pt}
      \begin{block}{And the surprise}
        $f(4x)=\dfrac{1}{4x}=\dfrac{1/4}{x}$ \\[3pt]
        On \emph{this} parent, $f(kx)$ and $kf(x)$ give the same family. One letter is enough.
      \end{block}
    \end{column}
  \end{columns}
\end{frame}

% GENERAL FORM
\begin{frame}[t, plain]
  \forestheader{All together: $g(x)=\dfrac{a}{x-h}+k$}
  \sectionlabel[goldacc]{General form}
  \begin{columns}[T]
    \begin{column}{0.42\textwidth}
      \centering
      \begin{tikzpicture}[scale=0.30]
        \lgrid{-3}{8}{-4}{6}
        \draw[navy, dashed, thick] (3,-4)--(3,6);
        \draw[navy, dashed, thick] (-3,1)--(8,1);
        \node[navy, font=\tiny, above] at (3,6) {$x=3$};
        \node[navy, font=\tiny, below left] at (-0.6,1) {$y=1$};
        \begin{scope} \clip (-3,-4) rectangle (8,6);
          \draw[very thick, forest, domain=-3:2.6, samples=100, smooth] plot (\x,{2/((\x)-3)+1});
          \draw[very thick, forest, domain=3.4:8, samples=100, smooth] plot (\x,{2/((\x)-3)+1});
        \end{scope}
        \fill[charcoal] (1,0) circle (3.4pt);
      \end{tikzpicture} \\[2pt]
      {\footnotesize $g(x)=\dfrac{2}{x-3}+1$}
    \end{column}
    \begin{column}{0.55\textwidth}
      \begin{itemize}\setlength{\itemsep}{3pt}
        \item $x-3=0 \Rightarrow$ \textbf{VA} $x=3$.
        \item The outside $+1$ \textbf{is} the \textbf{HA} $y=1$.
        \item \textbf{Center} $(3,1)$ --- where they cross, and \emph{not on the graph}.
        \item Domain $x\neq3$; \; Range $y\neq1$.
        \item $y$-int: $g(0)=\frac{2}{-3}+1=\frac13$. \; $x$-int: $\frac{2}{x-3}=-1 \Rightarrow (1,0)$.
      \end{itemize}
      \vspace{3pt}
      {\footnotesize Intercepts are \textbf{computed}, never read off the picture.}
    \end{column}
  \end{columns}
\end{frame}

% NEVER CROSSES
\begin{frame}[t, plain]
  \forestheader{Does it ever cross its horizontal asymptote?}
  \sectionlabel[goldacc]{Stronger than 5.0}
  \begin{columns}[T]
    \begin{column}{0.50\textwidth}
      \[ g(x)=1 \quad\Longrightarrow\quad \frac{2}{x-3}=0 \]
      \vspace{-6pt}
      \begin{itemize}\setlength{\itemsep}{3pt}
        \item A fraction is zero only when its \textbf{numerator} is zero.
        \item The numerator is $2$.
        \item \textbf{Impossible.} No solution, ever.
      \end{itemize}
    \end{column}
    \begin{column}{0.46\textwidth}
      \begin{block}{But be careful what you conclude}
        A \emph{transformed parent} never crosses its HA --- its numerator is a nonzero constant.
        \\[3pt]
        Rational functions in general \textbf{may} cross: $\frac{2x}{x^2+1}$ crosses $y=0$ at the
        origin (Lesson 5.0).
      \end{block}
    \end{column}
  \end{columns}
\end{frame}

% GRAPH TO EQUATION
\begin{frame}[t, plain]
  \forestheader{Backwards: write the equation from the graph}
  \sectionlabel[goldacc]{A2.F.1b}
  \begin{columns}[T]
    \begin{column}{0.40\textwidth}
      \centering
      \begin{tikzpicture}[scale=0.30]
        \lgrid{-6}{5}{-6}{3}
        \draw[navy, dashed, thick] (-1,-6)--(-1,3);
        \draw[navy, dashed, thick] (-6,-2)--(5,-2);
        \node[navy, font=\tiny, above right] at (-1,2.2) {$x=-1$};
        \node[navy, font=\tiny, above left] at (4.6,-2) {$y=-2$};
        \begin{scope} \clip (-6,-6) rectangle (5,3);
          \draw[very thick, forest, domain=-6:-1.35, samples=100, smooth] plot (\x,{2/((\x)+1)-2});
          \draw[very thick, forest, domain=-0.6:5, samples=100, smooth] plot (\x,{2/((\x)+1)-2});
        \end{scope}
        \fill[charcoal] (0,0) circle (3.4pt);
        \fill[charcoal] (1,-1) circle (3.4pt) node[above right, font=\tiny]{$(1,-1)$};
      \end{tikzpicture}
    \end{column}
    \begin{column}{0.57\textwidth}
      \begin{enumerate}\setlength{\itemsep}{2pt}
        \item VA $x=-1 \Rightarrow h=-1$.
        \item HA $y=-2 \Rightarrow k=-2$.
        \item $y=\dfrac{a}{x+1}-2$.
        \item Substitute $(1,-1)$: \; $-1=\dfrac a2-2 \Rightarrow a=2$.
        \item \textbf{Check on a second point}: at $(0,0)$, $\frac21-2=0$ \checkmark
      \end{enumerate}
      \vspace{3pt}
      \begin{block}{Watch the sign}
        The wall is at $x=-1$, and the equation shows $x+1$, because $x-(-1)=x+1$.
      \end{block}
    \end{column}
  \end{columns}
\end{frame}

% COMPARE + DISGUISE
\begin{frame}[t, plain]
  \forestheader{Compare, contrast --- and one disguise}
  \sectionlabel{A2.F.1a/e}
  \begin{columns}[T]
    \begin{column}{0.50\textwidth}
      \centering
      \renewcommand{\arraystretch}{1.2}
      {\footnotesize
      \begin{tabular}{|l|c|c|}
      \hline
       & $y=x^2$ or $y=|x|$ & $y=\frac1x$ \\ \hline
      Domain     & all reals & $x\neq0$ \\ \hline
      Range      & $y\ge0$   & $y\neq0$ \\ \hline
      Asymptotes & none      & $x=0$, $y=0$ \\ \hline
      Intercepts & $(0,0)$   & none \\ \hline
      Pieces     & one       & two branches \\ \hline
      Symmetry   & $y$-axis  & origin \\ \hline
      \end{tabular}}
    \end{column}
    \begin{column}{0.46\textwidth}
      \begin{block}{The disguise}
        \[ \frac{x-1}{x-3}=\frac{(x-3)+2}{x-3}=1+\frac{2}{x-3} \]
        It was the anchor all along. VA $x=3$, HA $y=1$.
      \end{block}
      \vspace{3pt}
      {\footnotesize Top and bottom both \textbf{linear} $\Rightarrow$ a shifted, stretched $\frac1x$.
      And the leftover $k$ \emph{is} the ratio of the leading coefficients --- which is why Lesson
      5.0's degree rule was true.}
    \end{column}
  \end{columns}
\end{frame}

% CLOSE
\begin{frame}[t, plain]
  \forestheader{Today in one line --- and where it goes next}
  \sectionlabel[goldacc]{Close}
  \begin{block}{}
    \centering\textbf{The asymptotes are the parent's axes, and they travel with the graph.} \\[3pt]
    $x=h$ from the denominator (which does the opposite of what it says); $y=k$ from the number
    outside; $a$ moves neither.
  \end{block}
  \vspace{6pt}
  \begin{columns}[T]
    \begin{column}{0.48\textwidth}
      \begin{block}{Errors that cost the most}
        \begin{itemize}\setlength{\itemsep}{2pt}
          \item Flipping the sign of $h$.
          \item Reading the HA off the numerator.
          \item Forgetting the \emph{range} exclusion.
          \item Plotting the center as a point of the graph.
        \end{itemize}
      \end{block}
    \end{column}
    \begin{column}{0.48\textwidth}
      \begin{block}{Next: Lesson 5.5}
        The denominator factors into \emph{two} pieces. The asymptotes stop being visible in the
        equation, and the graph gets assembled instead of read --- holes vs.\ walls, a degree
        comparison, intercepts, and a sign chart.
      \end{block}
    \end{column}
  \end{columns}
\end{frame}

\end{document}
